\chapter{Literature Review}

\section{Related Work}
\TemplateTodo{Summarize the most relevant related studies. Explain their objectives, methods, data, results, and limitations. End by positioning your study against those works.}

A literature review is not merely a list of article summaries. Organize it thematically so readers understand the development of knowledge, remaining gaps, and why this study is needed.

A summary of the directly comparable studies is presented in Table \ref{tab:related-work}. Mention a table in a sentence before it appears, because readers need to know what they are about to read.

% TABLE EXAMPLE 1: every column flexible, using L{...} weights.
%
% A related-work table is usually long, so it uses xltabular: the width adapts
% like tabularx, but the body may break across pages like longtable. That is why
% the table is NOT wrapped in a table environment and takes no [H]: xltabular
% places itself.
%
% The header rows between \endfirsthead and \endhead repeat on every
% continuation page, so readers never have to flip back to learn what a column
% means. \endlastfoot closes the table with one \hline, so the final data row
% needs no \hline of its own.
%
% L{weight} is a weighted, ragged-right X column. The weights in one table must
% sum to the number of columns: the five below add up to
% 0.70 + 1.05 + 1.20 + 1.15 + 0.90 = 5.00. Give the wordiest column the largest
% weight. The definition lives in config/thesis-preamble.tex.
\begin{xltabular}{\linewidth}{|L{0.70}|L{1.05}|L{1.20}|L{1.15}|L{0.90}|}
\caption{Summary of directly comparable studies.}\label{tab:related-work}\\
\hline
Author (Year) & Method & Dataset & Results & Limitations \\
\hline
\endfirsthead
\hline
Author (Year) & Method & Dataset & Results & Limitations \\
\hline
\endhead
\hline
\endlastfoot
Author name (year) \cite{exampleReference} & Short summary of the method used & Dataset name, size, and coverage & The main reported figures & The gap your study addresses \\
\hline
Author name (year) & Add one row per study you review & Name the data source & State results with clear units & Write the limitation specifically \\
\end{xltabular}

\section{Theoretical Foundation}
\TemplateTodo{Explain concepts, theories, algorithms, standards, or frameworks that are actually used in the study.}

\subsection{First Main Concept}
Write the definition, working principle, and relevance of the concept to the study. Use citations for definitions or theories that are not your own opinion.

\subsection{Second Main Concept}
Add other concepts as needed. If using formulas, explain the meaning of each symbol and why the formula is used.

\section{Conceptual Framework}
\TemplateTodo{Explain the relationship between the problem, theory, method, data, and research outputs. A conceptual diagram can help this section.}

\begin{figure}[H]
\centering
\includegraphics[width=0.75\linewidth]{example-figure.pdf}
\caption{Example figure usage in the template}
\label{fig:example-figure}
\end{figure}

An example of figure usage is shown in Figure~\ref{fig:example-figure}. Mathematical notation can be used to explain the relationship between layers in Figure~\ref{fig:example-figure}. Let the input vector be denoted as
\begin{equation}
\mathbf{x} =
\begin{bmatrix}
x_1 & x_2 & \cdots & x_n
\end{bmatrix}^{\mathsf{T}}.
\end{equation}
For the $l$-th hidden layer, neuron activations are computed by
\begin{equation}
\mathbf{a}^{[l]} = \sigma\left(\mathbf{W}^{[l]}\mathbf{a}^{[l-1]} + \mathbf{b}^{[l]}\right),
\qquad l = 1,2,\ldots,L-1,
\end{equation}
where $\mathbf{a}^{[0]}=\mathbf{x}$, $\mathbf{W}^{[l]}$
is the weight matrix at layer $l$,
$\mathbf{b}^{[l]}$ is the bias vector, and
$\sigma(\cdot)$ is the activation function. The model output is denoted as
\begin{equation}
\hat{\mathbf{y}} = f\left(\mathbf{W}^{[L]}\mathbf{a}^{[L-1]} + \mathbf{b}^{[L]}\right),
\end{equation}
with output components
\begin{equation}
\hat{\mathbf{y}} =
\begin{bmatrix}
\hat{y}_1 & \hat{y}_2 & \cdots & \hat{y}_k
\end{bmatrix}^{\mathsf{T}}.
\end{equation}

\subsection{Figures With Panel Labels}
A figure made of several panels is labelled (a), (b), and so on. Instead of several \texttt{figure} environments that would each take their own number, one image is wrapped in a \texttt{tikzpicture} and the panel labels are placed as nodes beneath it, so all panels keep sharing a single number and a single caption.

\begin{figure}[H]
\centering
% The ($A!0.25!B$) syntax needs the calc library and means "the point 25% of the
% way from A to B", so a panel label always sits directly under the centre of
% its panel no matter how wide the image is scaled.
\begin{tikzpicture}
  \node[anchor=south west, inner sep=0pt] (img) at (0,0)
    {\includegraphics[width=0.85\linewidth]{example-figure.pdf}};
  \node[anchor=north, font=\small, yshift=-2mm]
    at ($(img.south west)!0.25!(img.south east)$) {(a) before processing};
  \node[anchor=north, font=\small, yshift=-2mm]
    at ($(img.south west)!0.75!(img.south east)$) {(b) after processing};
\end{tikzpicture}
\caption{Example figure with two panel labels}
\label{fig:example-panels}
\end{figure}

The two conditions are compared in Figure~\ref{fig:example-panels}. For figures drawn entirely in TikZ or pgfplots, work in the \texttt{tikz/} folder and include the resulting PDF; see \texttt{tikz/README.md}.
